\section{Biology for machine-learning readers}
\label{sec:bio}

\subsection{Proteolysis and the substrate window}
A protease is an enzyme that hydrolyses a peptide bond inside a protein or peptide. The bond that is cut is the \emph{scissile bond}. Residues on its amino-terminal side are numbered outward from the bond as $\pos{1},\pos{2},\dots$ and residues on the carboxy-terminal side as $\prm{1},\prm{2},\dots$ (in Kridel et al.\ 2001 the scissile bond is drawn as ``$\downarrow$''; residues to its right are the primed residues). The enzyme surface that receives residue $\pos{i}$ is called subsite $\mathrm{S}i$, and $\prm{i}$ binds $\mathrm{S}i'$. Specificity is therefore \emph{positional}: a protease may want proline at $\pos{3}$ and a large hydrophobic residue at $\prm{1}$, and those preferences are largely independent of one another. This is why an additive, position-specific model is a natural starting point.

The unit of data throughout this report is the eight-residue window
\[
\underbrace{\pos{4}\ \pos{3}\ \pos{2}\ \pos{1}}_{\text{non-primed}}\ \ \downarrow\ \ \underbrace{\prm{1}\ \prm{2}\ \prm{3}\ \prm{4}}_{\text{primed}},
\]
so the scissile bond always sits between the fourth and fifth character. As a worked example, the MMP-9 substrate SGKIPRT$\downarrow$LTA measured by Kridel et al.\ (2001) has $\pos{1}=\mathrm{T}$, $\pos{2}=\mathrm{R}$, $\pos{3}=\mathrm{P}$, $\pos{4}=\mathrm{I}$ and $\prm{1}=\mathrm{L}$, $\prm{2}=\mathrm{T}$, $\prm{3}=\mathrm{A}$; its window is \texttt{IPRTLTA}$-$ with $\prm{4}$ unknown because the peptide ends there.

\begin{bionote}[Why the alignment matters so much]
Because specificity is positional, shifting a peptide by one residue changes which enzyme subsite every residue is tested against. A correct sequence at the wrong offset is, to the model, a different and usually poor substrate. Section~\ref{sec:design} shows a real instance of this error in the sensor design and its consequences.
\end{bionote}

\subsection{The proteases}
\begin{table}[!htbp]
\centering\small
\begin{tabular}{@{}>{\raggedright\arraybackslash}p{5.4cm}l>{\raggedright\arraybackslash}p{3.0cm}>{\raggedright\arraybackslash}p{4.4cm}@{}}\toprule
Protease & MEROPS id & Class & Structure used here \\\midrule
MMP-2 (matrix metallopeptidase-2) & M10.003 & zinc metalloprotease & 1QIB (catalytic domain) \\
MMP-8 (matrix metallopeptidase-8) & M10.002 & zinc metalloprotease & 1JAN (with Pro-Leu-Gly-hydroxylamine) \\
MMP-9 (matrix metallopeptidase-9) & M10.004 & zinc metalloprotease & 1GKC (with an inhibitor) \\
MMP-14 (membrane-type MMP-1) & M10.014 & zinc metalloprotease & not shown \\
Neutrophil elastase (elastase-2) & S01.131 & serine protease & 1HNE (with MeOSuc-AAPA-CMK) \\
LasB (pseudolysin, \emph{P.\ aeruginosa}) & M04.005 & zinc metalloprotease & 1EZM \\
Faecal group (trypsin 1, chymotrypsin B, elastase-1) & S01.151--153 & serine proteases & not shown \\\bottomrule
\end{tabular}
\caption{Panel proteases. Names and identifiers are taken from the MEROPS records used to extract the training data; the ``faecal'' group pools three unrelated digestive serine proteases from different species (bovine, human, porcine), which the data labels together.}
\label{tab:panel}
\end{table}

\paragraph{Matrix metalloproteinases (MMPs).} MMPs are zinc-dependent endopeptidases. Their catalytic domain of about 170 residues contains the zinc-binding motif HEXXHXXGXXH and a conserved methionine forming a ``Met-turn''; the catalytic domains also carry a structural zinc ion and two to three calcium ions (Nagase and Woessner, 1999). MMP-2 and MMP-9 have three fibronectin type II repeats inserted into the catalytic domain (Nagase and Woessner, 1999). We recover the histidine ligands of the catalytic zinc directly from the crystal coordinates in Section~\ref{sec:structures}.

\paragraph{MMP-9 specificity.} Kridel et al.\ (2001) selected MMP-9 substrates by phage display and measured kinetics for representative peptides. They report that most selected substrates contain a hydrophobic residue at $\prm{1}$, consistent with the deep hydrophobic S1$'$ pocket of MMPs; that the best (group~I) substrates carry Pro at $\pos{3}$ together with a large hydrophobic residue at $\prm{1}$; that Arg at $\pos{2}$ (and often $\pos{1}$) is favoured; and that Ser/Thr is favoured at $\prm{2}$. They derive the consensus Pro-Arg-Ser/Thr$\downarrow$Hy-Ser/Thr spanning $\pos{3}$ to $\prm{2}$. They also note that ``many different MMPs recognize the Pro-X-X-Hy core sequence'' and that ``substrate selectivity can be conferred by subsite interactions outside of the dominant P3 and P1$'$ subsites''.

\begin{bionote}[Consequence for selectivity]
If the dominant preferences (Pro at $\pos{3}$, hydrophobic at $\prm{1}$) are shared across the MMP family, a peptide that is an excellent substrate for one MMP will usually be a good substrate for several. Separating them relies on the weaker, secondary subsite preferences, which is a statistically harder problem. We will see this quantitatively in Table~\ref{tab:selmat}.
\end{bionote}

\paragraph{Elastases, LasB and digestive proteases.} Neutrophil elastase (MEROPS S01.131, entry ELNE\_HUMAN) is a serine protease released by neutrophils; in the crystal structure used here its catalytic triad is His57, Asp102 and Ser195. LasB (pseudolysin, MEROPS M04.005) is a zinc metalloprotease of \emph{Pseudomonas aeruginosa}; Yang et al.\ (2015) studied its degradation of elastin. The faecal group is a pooled set of digestive serine proteases; in our data it is dominated by a Lys/Arg preference at $\pos{1}$ (Figure~\ref{fig:logos}), and it serves mainly as a strongly distinct comparison class.

\subsection{Wound context}
Trengove et al.\ (1999) summarise that MMPs are present in elevated amounts in wound fluid from chronic leg ulcers compared with acute wounds, and cite work reporting MMP-9 (and to a lesser extent MMP-2) elevated ten- to twenty-five-fold in decubitus patients compared with mastectomy fluid, and MMP-9 increased with TIMPs decreased in chronic wounds. They also cite reports that neutrophil elastase degrades fibronectin in chronic wounds and that its natural inhibitor $\alpha_1$-antitrypsin is low, while noting that other reports found low elastase. The evidence for elastase is therefore mixed in the source we hold. This is background only: none of our models use clinical data.

\subsection{How each data source measures ``cleavage''}
\paragraph{MEROPS.} A curated database of cleavage sites collected from the literature, listing for each site the eight residues around the scissile bond and the peptidase. It mixes physiological and synthetic substrates. \paragraph{PICS (Eckhard et al., 2016).} Proteomic identification of cleavage sites incubates an MMP with peptide libraries derived from a human proteome digested with trypsin or with GluC; cleavage products are identified by LC-MS/MS (Mascot, X!\,Tandem, iProphet, 1\% false-discovery rate), and every reported peptide carries a chemically labelled N-terminus (thioacyl, $+89.01$~Da, on $81\%$ of the peptides we use; the remaining $19\%$ carry a $+29.04$~Da label that the source legend does not define). The non-primed residues $\pos{4}$--$\pos{1}$ are recovered by locating the peptide in its source protein sequence. Because the libraries are generated with trypsin or GluC, their composition is constrained at peptide ends, which can leak into primed positions of short peptides; we do not correct for this. \paragraph{Yang et al.\ (2015).} Peptides released from bovine elastin by pseudolysin were identified by mass spectrometry; both termini of each released peptide are cleavage sites. The source is a single protein, which matters in Section~\ref{sec:limits}.
